\section{Ряды}

\subsection{Простейшие свойства}

\begin{defn}{Основные понятия}{}
    \( \{ a_n \}_{n = 1}^{\infty} \)

    \( S_n = \sum\limits_{k = 1}^{n} a_k \) --- частичные суммы.

    \( \sum\limits_{k = 1}^{\infty} a_k \) называется сходящимся,
    если сходится \( \{ S_n \}_{n = 1}^{\infty} \)
    и
    \( \sum\limits_{k = 1}^{\infty} a_k = \lim\limits_{n \to \infty} S_n \)
\end{defn}

\begin{example}{}{}
    \( \sum\limits_{n = 0}^{\infty} q^n \) сходится
    тогда и только тогда, когда \( |q| < 1 \).

    Сумма вычисляется по формуле \( \sum\limits_{n = 0}^{\infty} q^n = \frac{1}{1 - q} \).
\end{example}

\begin{assrt}{Критерий Коши}{}
    \( \sum\limits_{n = 1}^{\infty} a_n \)
    сходится
    \( \Leftrightarrow \)
    \(
        \forall \varepsilon > 0 \ \exists N \in \NN :
        \forall n > N, \ p \in \NN \
        \left| \sum\limits_{k = n + 1}^{n + p} a_k \right| < \varepsilon
    \)
\end{assrt}

Вспомним несколько простых утверждений:
\begin{enumerate}
    \item
        Если \( \sum\limits_{n = 1}^{\infty} a_n \) и \( \sum\limits_{n = 1}^{\infty} b_n \)
        сходятся, то
        \(
            \sum\limits_{n = 1}^{\infty} (\alpha a_n + \beta b_n)
            =
            \alpha \sum\limits_{n = 1}^{\infty} a_n
            +
            \beta \sum\limits_{n = 1}^{\infty} b_n
        \)
    \item
        \( \sum\limits_{n = 1}^{\infty} a_n \) сходится
        \( \Leftrightarrow \)
        \( \sum\limits_{n = p}^{\infty} a_n \) сходится.
\end{enumerate}

\begin{assrt}{Признак сравнения}{}
    \( \{ a_n \geq 0 \}_{n = 1}^{\infty} \) \quad \( \{ b_n \geq 0 \}_{n = 1}^{\infty} \)

    \( a_n \leq b_n \ \forall n \) и \( \sum\limits_{n = 1}^{\infty} b_n \) сходится
    \( \Rightarrow \)
    \( \sum\limits_{n = 1}^{\infty} a_n \) сходится.
\end{assrt}

\textit{Доказательство:}

\[
    \left| \sum\limits_{k = N + 1}^{N + p} a_k \right| \leq \left| \sum\limits_{k = N + 1}^{N + p} b_k \right|
\]

Ура, критерий Коши. \qed

\subsection{Неотрицательные ряды}

\begin{assrt}{Признак Коши}{}
    Рассмотрим \( \{ a_n \} \), где \( a_n \geq 0 \ \forall n \).
    \begin{enumerate}
        \item
            Пусть \( \exists \delta > 0 \) такое, что \( \sqrt[n]{a_n} \leq 1 - \delta \), начиная
            с некоторого \( n \). Тогда \( \sum\limits_{n = 1}^{\infty} a_n \) сходится.
        \item
            Пусть \( \sqrt[n]{a_n} \geq 1 \), начиная с некоторого \( n \).
            Тогда \( \sum\limits_{n = 1}^{\infty} a_n \) расходится.
    \end{enumerate}
\end{assrt}

\textit{Доказательство:}
\begin{enumerate}
    \item
        \( \sqrt[n]{a_n} \leq 1 - \delta \Rightarrow a_n \leq (1 - \delta)^n \Rightarrow \)
        сходится по признаку сравнения.
    \item
        \( \sqrt[n]{a_n} \geq 1 \Rightarrow a_n \geq 1 \)
\qed
\end{enumerate}

\begin{assrt}{Интегральный признак}{}
    \( f : [1, \infty) \to \RR_{\geq 0} \) интегрируема по Риману и монотонно невозрастает.

    Тогда \( \sum\limits_{n = 1}^{\infty} f(n) \) сходится
    \( \Leftrightarrow \)
    \( \int\limits_{1}^{\infty} f(x) dx \) сходится.
\end{assrt}

\textit{Доказательство:}
\[
    S_k
    =
    \sum\limits_{n = 1}^{k} a_n
    =
    \sum\limits_{n = 1}^{k} f(n)
    \geq
    \int\limits_{1}^{k + 1} f(x) dx
    \geq
    \sum\limits_{n = 2}^{k + 1} f(n)
    =
    \sum\limits_{n = 2}^{k + 1} a_n
    =
    S_{k + 1} - a_1
\]

Получили неравенство в обе стороны, а значит доказали. \qed

\begin{example}{}{}
    Ряд \( \sum\limits_{n = 1}^{\infty} \frac{1}{n^a} \) сходится
    \( \Leftrightarrow \)
    \( \int\limits_{1}^{\infty} \frac{1}{x^a} dx \) сходится,
    то есть при \( a > 1 \).
\end{example}

\begin{thrm}{Признак Куммера}{}
    \( \{ a_n > 0 \}_{n = 1}^{\infty} \), \( \{ c_n > 0 \}_{n = 1}^{\infty} \)

    \begin{enumerate}
        \item
            \( \exists \delta > 0 : \frac{a_n}{a_{n + 1}} c_n - c_{n + 1} > \delta \)
            начиная с некоторого \( n \)
            \( \Rightarrow \)
            \( \sum\limits_{n = 1}^{\infty} a_n \) сходится.
        \item
            \( \sum\limits_{n = 1}^{\infty} \frac{1}{c_n} \) расходится
            и \( \frac{a_n}{a_{n + 1}} c_n - c_{n + 1} \leq 0 \) с некоторого \( n \)
            \( \Rightarrow \) \( \sum\limits_{n = 1}^{\infty} a_n \) расходится.
    \end{enumerate}
\end{thrm}

\begin{example}{}{}
    \begin{enumerate}
        \item
            \( c_n = 1 \) --- признак Даламбера

            \( \frac{a_n}{a_{n + 1}} > 1 + \delta \) начиная с некоторого \( n \)
            \( \Rightarrow \)
            \( \sum\limits_{n = 1}^{\infty} a_n \) сходится.

            \( \frac{a_n}{a_{n + 1}} \leq 1 \) начиная с некоторого \( n \)
            \( \Rightarrow \)
            \( \sum\limits_{n = 1}^{\infty} a_n \) расходится.
        \item
            \( c_n = n \) --- признак Раабе

            \( \frac{a_n}{a_{n + 1}} > 1 + \frac{1 + \delta}{n} \) начиная с некоторого \( n \)
            \( \Rightarrow \)
            \( \sum\limits_{n = 1}^{\infty} a_n \) сходится.

            \( \frac{a_n}{a_{n + 1}} \leq 1 + \frac{1}{n} \) начиная с некоторого \( n \)
            \( \Rightarrow \)
            \( \sum\limits_{n = 1}^{\infty} a_n \) расходится.
    \end{enumerate}
\end{example}

\textit{Доказательство:}

\begin{enumerate}
    \item
        Будем считать, что условие выполняется с \( n = 1 \)

        \( a_n c_n - a_{n + 1} c_{n + 1} > \delta a_{n + 1} > 0 \)

        Последовательность \( \{ a_n c_n \}_{n = 1}^{\infty} \) убывает
        \( \Rightarrow \)
        \( \{ a_n c_n \}_{n = 1}^{\infty} \) сходится.

        Значит \( \sum\limits_{n = 1}^{\infty} ( a_n c_n - a_{n + 1} c_{n + 1} ) \) сходится
        \( \Rightarrow \)
        \( \sum\limits_{n = 1}^{\infty} \delta a_{n + 1} \) сходится
        \( \Rightarrow \)
        \( \sum\limits_{n = 1}^{\infty} a_n \) сходится.
    \item
        Введем \( d_n = \frac{1}{c_n} \).

        Тогда
        \( \frac{a_n}{a_{n + 1}} \leq \frac{c_{n + 1}}{c_n} = \frac{d_n}{d_{n + 1}} \).

        Получаем:
        \begin{gather*}
            \frac{a_n}{a_{n + 1}} \cdot \frac{a_{n + 1}}{a_{n + 2}} \cdot \ldots \cdot \frac{a_m}{a_{m + 1}}
            \leq
            \frac{d_n}{d_{n + 1}} \cdot \frac{d_{n + 1}}{d_{n + 2}} \cdot \ldots \cdot \frac{d_m}{d_{m + 1}}
            \\
            \frac{a_n}{a_{m + 1}} \leq \frac{d_n}{d_{m + 1}}
            \\
            a_{m + 1} \geq d_{m + 1} \frac{a_n}{d_n}
        \end{gather*}

        Давайте зафиксируем \( n \): возьмем любое, для которого выполняются
        условия признака, тогда \( \frac{a_n}{d_n} \) --- какая-то константа.
        Но \( \sum\limits_{n = 1}^{\infty} d_n = \sum\limits_{n = 1}^{\infty} \frac{1}{c_n} \)
        расходится, а значит расходится и \( \sum\limits_{n = 1}^{\infty} a_n \).
\qed
\end{enumerate}

\begin{lemma}{}{}
    \( \{ a_n \geq 0 \}_{n = 1}^{\infty} \)

    Ряд \( \sum\limits_{n = 1}^{\infty} a_n \) сходится
    \( \Leftrightarrow \)
    \( \sup \left\{ \sum\limits_{n \in F} a_n : \underset{\text{конечно}}{F \subset \NN} \right\} < +\infty \)

    Более того,
    \(
        \sum\limits_{n = 1}^{\infty} a_n
        =
        \sup \left\{ \sum\limits_{n \in F} a_n : \underset{\text{конечно}}{F \subset \NN} \right\}
    \)
\end{lemma}

\textit{Доказательство:}

\begin{itemize}
    \item
        \( \Rightarrow \)

        \( F \subset \NN \) --- конечное \( \Rightarrow \exists N : F \subset \{ 1, 2, \ldots, N \} \)
        \( \Rightarrow \)
        \(
            \sum\limits_{n \in F}^{} a_n
            \leq
            \sum\limits_{n = 1}^{N} a_n
            \leq
            \sum\limits_{n = 1}^{\infty} a_n
        \)

        Из этого следует, что
        \(
            \sup \left\{ \sum\limits_{n \in F} a_n : \underset{\text{конечно}}{F \subset \NN} \right\}
            \leq
            \sum\limits_{n = 1}^{\infty} a_n
        \)
    \item
        \( \Leftarrow \)

        \(
            S_N
            =
            \sum\limits_{n = 1}^{N} a_n
            \leq
            \sum\limits_{ n \in \{ 1, 2, \ldots, N \} }^{} a_n
            \leq
            \sup \left\{ \sum\limits_{n \in F} a_n : \underset{\text{конечно}}{F \subset \NN} \right\}
        \)

        Отсюда явно следует, что \( \{ S_n \} \) сходится, причем
        \[
            \sum\limits_{n = 1}^{\infty} a_n
            =
            \lim\limits_{n \to \infty} S_n
            \leq
            \sup \left\{ \sum\limits_{n \in F} a_n : \underset{\text{конечно}}{F \subset \NN} \right\}
        \]
\qed
\end{itemize}

\begin{coroll}{}{}
    \( \{ a_n \geq 0 \}_{n = 1}^{\infty} \) и \( \tau : \NN \to \NN \) --- биекция.

    Тогда \( \sum\limits_{n = 1}^{\infty} a_n \) сходится
    \( \Leftrightarrow \)
    \( \sum\limits_{n = 1}^{\infty} a_{\tau(n)} \) сходится.

    Причем, если ряды сходятся, то их суммы равны.
\end{coroll}

\begin{coroll}{}{}
    Пусть \( \{ A_n \subset \NN \}_{n = 1}^{\infty} \)
    такие, что \( A_n \cap A_m = \varnothing \ (n \neq m) \)
    и \( \NN = \bigcup\limits_{n = 1}^{\infty} A_n \).

    \( \{ a_m \geq 0 \}_{m = 1}^{\infty} \)

    Тогда
    \(
        \sum\limits_{m = 1}^{\infty} a_m
        =
        \sum\limits_{n = 1}^{\infty} \left( \sum\limits_{m \in A_n} a_m \right)
    \)
\end{coroll}

\newpage

\textit{Доказательство:}

\begin{itemize}
    \item
        \( \leq \)

        \(
            \sum\limits_{m = 1}^{\infty} a_m
            =
            \sup \left\{ \sum\limits_{n \in F} a_n : \underset{\text{конечно}}{F \subset \NN} \right\}
        \)

        \( F \subset \NN \) --- конечное,
        \(
            F \subset \bigcup\limits_{n = 1}^{\infty} A_n
            \Rightarrow
            \exists N \in \NN : F \subset \bigcup\limits_{n = 1}^{N} A_n
        \)

        \[
            \sum\limits_{m \in F} a_m
            =
            \sum\limits_{n = 1}^{N} \left( \sum\limits_{m \in F \cap A_n} a_m \right)
            \leq
            \sum\limits_{n = 1}^{N} \left( \sum\limits_{m \in A_n} a_m \right)
            \leq
            \sum\limits_{n = 1}^{\infty} \left( \sum\limits_{m \in A_n} a_m \right)
        \]
    \item
        \( \geq \)

        Пусть \( \sum\limits_{m = 1}^{\infty} a_m < +\infty \)
        (иначе бесконечность \( \geq \) чего угодно).

        Давайте докажем, что
        \(
            \sum\limits_{n = 1}^{N} \left( \sum\limits_{m \in A_n}^{} a_m \right)
            \leq
            \sum\limits_{m = 1}^{\infty} a_m
        \)
        для всех \( N \in \NN \).

        Зафиксируем \( \varepsilon > 0 \)
        и для \( n \) от \( 1 \) до \( N \)
        выберем конечное \( F_n \subseteq A_n \) так, что
        \[
            \sum\limits_{m \in F_n}^{} a_m \geq \sum\limits_{m \in A_n}^{} a_m - \frac{\varepsilon}{N}
        \]

        Просуммируем:
        \[
            \sum\limits_{n = 1}^{N} \left( \sum\limits_{m \in A_n}^{} a_m \right)
            \leq
            \sum\limits_{n = 1}^{N} \left( \sum\limits_{m \in F_n}^{} a_m + \frac{\varepsilon}{N} \right)
            =
            \varepsilon
            +
            \sum\limits_{m \in F_1 \cup \ldots \cup F_N}^{} a_m
            \leq
            \varepsilon
            +
            \sum\limits_{m = 1}^{\infty} a_m
        \]

        Поскольку \( \varepsilon \) был произвольным, доказали то, что хотели.

        Остается сделать предельный переход по \( N \).
\qed
\end{itemize}

\subsection{Знакопеременные ряды}

\begin{defn}{}{}
    Ряд \( \sum\limits_{n = 1}^{\infty} a_n \) называется сходящимся абсолютно,
    если \( \sum\limits_{n = 1}^{\infty} |a_n| < +\infty \).
    Если ряд сходится, но не абсолютно, он называется условно сходящимся.
\end{defn}

\begin{lemma}{}{}
    \( \sum\limits_{n = 1}^{\infty} a_n \)

    \( a_n^+ = \max \{ a_n, 0 \} \)

    \( a_n^- = \max \{ -a_n, 0 \} \)

    \( a_n = a_n^+ - a_n^- \)

    \begin{enumerate}
        \item
            Ряд \( \sum\limits_{n = 1}^{\infty} a_n \) сходится абсолютно
            \( \Leftrightarrow \)
            \( \sum\limits_{n = 1}^{\infty}  a_n^+ < +\infty \) и  \( \sum\limits_{n = 1}^{\infty} a_n^- < +\infty \)
        \item
            Если ряд \( \sum\limits_{n = 1}^{\infty} a_n \) сходится условно,
            то \( \sum\limits_{n = 1}^{\infty} a_n^+ = +\infty \)
            и \( \sum\limits_{n = 1}^{\infty} a_n^- = +\infty \)
    \end{enumerate}
\end{lemma}

\textit{Доказательство:}
\begin{enumerate}
    \item
        \begin{itemize}
            \item
                \( \Rightarrow \)

                \( a_n^+, a_n^- \leq |a_n| \) \( \Rightarrow \) признак сравнения
            \item
                \( \Leftarrow \)

                Следует из \( |a_n| = a_n^+ + a_n^- \)
        \end{itemize}
    \item
        \( a_n = a_n^+ - a_n^- \)

        Если два из трех рядов \( \sum\limits_{n = 1}^{\infty} a_n \),
        \( \sum\limits_{n = 1}^{\infty} a_n^+ \) и \( \sum\limits_{n = 1}^{\infty} a_n^- \)
        сходятся, то и третий тоже.
        Отсюда явно следует искомое.
\qed
\end{enumerate}

\begin{assrt}{}{}
    Если \( \sum\limits_{n = 1}^{\infty} a_n \) сходится абсолютно,
    \( \tau : \NN \to \NN \) --- биекция,
    то \( \sum\limits_{n = 1}^{\infty}  a_{\tau(n)} \)
    тоже сходится абсолютно, причем \( \sum\limits_{n = 1}^{\infty} a_n = \sum\limits_{n = 1}^{\infty} a_{\tau(n)} \).
\end{assrt}

\textit{Доказательство:}

Снова разобьем ряд на разность двух положительных:
\begin{gather*}
    a_n = a_n^+ - a_n^-
    \\
    a_{\tau(n)} = a_{\tau(n)}^+ - a_{\tau(n)}^-
    \\
    \sum\limits_{n = 1}^{\infty} a_n
    =
    \sum\limits_{n = 1}^{\infty} a_n^+
    -
    \sum\limits_{n = 1}^{\infty} a_n^-
    =
    \sum\limits_{n = 1}^{\infty} a_{\tau(n)}^+
    -
    \sum\limits_{n = 1}^{\infty} a_{\tau(n)}^-
    =
    \sum\limits_{n = 1}^{\infty} a_{\tau(n)}
\end{gather*}

Можно рассмотреть еще одно утверждение из неотрицательных рядов:
\begin{assrt}{}{}
    \( \sum\limits_{n = 1}^{\infty} a_n \) сходится абсолютно,
    \( \NN = \bigcup\limits_{n = 1}^{\infty} A_n \),
    \( \underset{n \neq m}{A_n \cap A_m} = \varnothing \).

    Тогда \( \sum\limits_{n = 1}^{\infty} a_n = \sum\limits_{n = 1}^{\infty} \left( \sum\limits_{m \in A_n}^{} a_m \right) \)
\end{assrt}

\begin{thrm}{Теорема Римана об условной сходимости}{}
    \( \sum\limits_{n = 1}^{\infty} a_n \) сходится условно
    \begin{enumerate}
        \item
            Тогда \( \forall s \in \RR \ \exists \tau : \NN \to \NN \) биекция такая,
            что \( \sum\limits_{n = 1}^{\infty} a_{\tau(n)} = s \)
        \item
            \( \exists \tau : \NN \to \NN \) биекция, такая что \( \sum\limits_{n = 1}^{\infty} a_{\tau(n)} \) расходится.
    \end{enumerate}
\end{thrm}

\textit{Доказательство (чуть-чуть рукомахание):}
\begin{enumerate}
    \item
        \( J^+ = \{ n \in \NN : a_n \geq 0 \} \)

        \( J^- = \{ n \in \NN : a_n < 0 \} \)

        Сначала набираем наименьшие индексы из \( J^+ \), пока сумма не больше \( s \):
        \[
            \sum\limits_{l = 1}^{k_1} a_{\tau(l)} > s
        \]

        Потом набираем наименьшие индексы из \( J^- \), пока сумма не меньше \( s \):
        \[
            \sum\limits_{l = 1}^{k_2} a_{\tau(l)} < s
        \]

        И так далее.

        В целом понятно, почему мы наберем все числа.

        Теперь рассмотрим частичную сумму \( S_k \).
        Пусть \( k_n \leq k \leq k_{n + 1} \).
        В таком случае выполнено:
        \[
            \left| \sum\limits_{l = 1}^{k} a_{\tau(l)} - s \right| < \max \{ | a_{\tau(k_n)} |, | a_{\tau(k_{n + 1})} | \}
        \]

        Поскольку изначальный ряд сходился, его члены стремились к нулю при увеличении \( n \),
        а значит и при перестановке это выполнится.
        Таким образом \( S_k \) стремится к \( s \).
    \item
        Давайте, как и в первом пункте, разделим ряд на положительные и отрицательные члены.
        Но положительными мы будем набирать сумму хотя бы \( 1 \),
        а отрицательными --- не больше чем \( -1 \).
        В таком случае, очевидно, полученный ряд разойдется.
\qed
\end{enumerate}

\begin{defn}{}{}
    \( \{ a_n \} \) имеет ограниченную вариацию,
    если \( \sum\limits_{n = 1}^{\infty} |a_{n + 1} - a_n| < +\infty \)
\end{defn}

\begin{example}{}{}
    \( \{ a_n \} \) монотонна и сходится \( \Rightarrow \) она имеет ограниченную вариацию.
\end{example}

\begin{assrt}{}{}
    \( \{ a_n \} \) имеет ограниченную вариацию
    \( \Leftrightarrow \)
    \( \exists \{ b_n \}, \{ c_n \} \) --- неубывающие и сходящиеся,
    такие что \( a_n = b_n - c_n \).
\end{assrt}

\textit{Доказательство:}

\begin{itemize}
    \item
        \( \Leftarrow \)

        Очевидно.
    \item
        \( \Rightarrow \)

        \( b_n = \sum\limits_{k = 1}^{n - 1} | a_{k + 1} - a_k | \), \( c_n = b_n - a_n \).

        Очевидно, что \( b_n \) неубывающая и сходящаяся.

        Покажем, что \( \{ c_n \} \) не убывает:
        \[
            c_{n + 1} - c_n = -a_{n + 1} + a_n + (b_{n + 1} - b_n)
            =
            -a_{n + 1} + a_n + | a_{n + 1} - a_n | \geq 0
        \]

        Заметим, что \( \{ a_n \} \) сходится.
        Действительно, вариация ограничена, а значит ряд сходится абсолютно,
        поэтому можно убрать модуль. А отсюда явно следует сходимость \( \{ a_n \} \).
        \( \{ b_n \} \) тоже сходится, откуда сходится и \( \{ c_n \} \).
\qed
\end{itemize}

\begin{thrm}{Признак Дирихле}{}
    \begin{enumerate}
        \item
            \( a_n \to 0 \) и имеет ограниченную вариацию
        \item
            \( \exists C : \forall k \in \NN \ \left| \sum\limits_{n = 1}^{k} b_n \right| \leq C \)
    \end{enumerate}

    Тогда ряд \( \sum\limits_{n = 1}^{\infty} a_n b_n \) сходится.
\end{thrm}

\textit{Доказательство:}

\( S_n = \sum\limits_{k = 1}^{n} b_k \), тогда:
\[
    \sum\limits_{k = 1}^{n} a_k b_k
    =
    \sum\limits_{k = 1}^{n} a_k (S_k - S_{k - 1})
    =
    -\sum\limits_{k = 1}^{n} (a_{k + 1} - a_k) S_k + a_{n + 1} S_n - a_1 S_0
\]

\( a_{n + 1} S_n \) стремится к нулю, \( a_1 S_0 \) является константой,
а \( \sum\limits_{k = 1}^{n} (a_{k + 1} - a_k) S_k \) сходится абсолютно.
А значит искомый ряд сходится.
\qed

\begin{assrt}{Признак Лейбница}{}
    Подставим \( b_n = (-1)^n \)

    \( \sum\limits_{n = 1}^{\infty} (-1)^n a_n \) сходится,
    если \( a_n \to 0 \) и имеет ограниченную вариацию (монотонна).
\end{assrt}

\begin{assrt}{Признак Абеля}{}
    \begin{enumerate}
        \item
            \( a_n \) имеет ограниченную вариацию (\( a_n \) не обязательно сходится к \( 0 \)).
        \item
            \( \sum\limits_{n = 1}^{\infty} b_n \) сходится.
    \end{enumerate}

    Тогда \( \sum\limits_{n = 1}^{\infty} a_n b_n \) сходится.
\end{assrt}

\subsection{Равномерная сходимость}

\( \{ f_n : X \to \RR \} \) --- последовательность функций.

\begin{defn}{}{}
    \begin{enumerate}
        \item
            Последовательность \( \{ f_n \} \) сходится поточечно,
            если для любого \( x \in X \) сходится \( \{ f_n (x) \} \).
        \item
            \( \{ f_n \} \) сходится равномерно,
            если она сходится поточечно и
            \[
                \forall \varepsilon > 0 \ \exists N \in \NN : | f_n (x) - f(x) | < \varepsilon \ \forall x \in X \ \forall n \geq N
            \]
        \item
            Ряд \( \sum\limits_{n = 1}^{\infty} f_n \) сходится поточечно (соответственно равномерно),
            если последовательность \( \left\{ \sum\limits_{k = 1}^{n} f_k \right\} \) сходится
            поточечно (равномерно).
    \end{enumerate}
\end{defn}

\begin{example}{}{}
    \( f_n (x) = \arctan (nx) \)

    \( \lim f_n (x) = \frac{\pi}{2} \sgn x \)
\end{example}

\begin{assrt}{}{}
    \( X \) --- метрическое пространство, \( f_n : X \to \RR \), \( x \in X \)

    Если \( f_n \) непрерывна в \( x \) для любого \( n \)
    и равномерно сходится к \( f \) в некоторой окрестности \( x \),
    то \( \lim f_n = f \) непрерывна в \( x \).
\end{assrt}

\textit{Доказательство:}

Хотим такое:
\[
    \forall \varepsilon > 0 \ \exists \delta : | f(y) - f(x) | < \varepsilon \text{ при } \rho(x, y) < \delta
\]
\begin{gather*}
    |f(y) - f(x)|
    =
    |f(y) - f_N(y) + f_N(y) - f_N(x) + f_N(x) - f(x)|
    \leq
    \\
    \leq
    |f(y) - f_N(y)| + |f_N(y) - f_N(x)| + |f_N(x) - f(x)|
\end{gather*}

Если \( f_n \to f \) равномерно на шаре \( B_\gamma(x) \), то
\begin{gather*}
    \exists N : | f_n(y) - f(y) | < \frac{\varepsilon}{3} \ \forall y \in B_\gamma (x) \ \forall n \geq N
    \\
    \exists \beta > 0 : | f_N (y) - f_N (x) | < \frac{\varepsilon}{3} \text{ при } \rho(y, x) < \beta
\end{gather*}

Положим \( \delta = \min \{ \gamma, \beta \} \), ура.
\qed

\begin{assrt}{}{}
    \( \{ f_n : [a, b] \to \RR \} \), \( f_n \to f \) равномерно, и \( f_n \) интегрируема
    по Риману для любого \( n \).
    Тогда \( f \) интегрируема по Риману и \( \int\limits_{a}^{b} f dx = \lim \int\limits_{a}^{b} f_n dx \)
\end{assrt}

\begin{example}{}{}
    Придумаем функцию на \( [0, 1] \)

    Пусть \( \{ q_n \}_{n = 1}^{\infty} \) --- все рациональные числа на отрезке \( [0, 1] \).

    \(
        f_n (x) = \begin{cases}
            1, \ x = q_j \text{ для некоторого } j \in \{ 1, 2, \ldots, n \}
            \\
            0, \text{ иначе}
        \end{cases}
    \)

    \( f = \lim f_n \) --- функция Дирихле --- не интегрируема по Риману.
\end{example}
